\def\stat{belozerov}





\def\tit{МАШИННОЕ ОБУЧЕНИЕ ДЛЯ ЗАДАЧ ЭКСТРАПОЛЯЦИИ С~МАЛЫМ ОБЪЕМОМ 


ДАННЫХ$^*$}





\def\titkol{Машинное обучение для задач экстраполяции с~малым объемом 


данных}





\def\aut{А.\,О.~Белозеров$^1$, А.\,И.~Мазур$^2$}





\def\autkol{А.\,О.~Белозеров, А.\,И.~Мазур}





\titel{\tit}{\aut}{\autkol}{\titkol}





\index{Белозеров А.\,О.}


\index{Мазур А.\,И.}


\index{Belozerov A.\,O.}


\index{Mazur A.\,I.}





{\renewcommand{\thefootnote}{\fnsymbol{footnote}}


\footnotetext[1]


{Статья публикуется по представлению программного комитета VI Международной на\-уч\-но-прак\-ти\-че\-ской 


конференции <<Информационные технологии и~высокопроизводительные вы\-чис\-ле\-ния>> (ITHPC-2021). 


Работа выполнена 


при поддержке Министерства науки и~высшего образования РФ в~рамках проекта №\,0818-2020-0005 и~гранта РФФИ  


№\,20-02-00357 (частично).}} 





\renewcommand{\thefootnote}{\arabic{footnote}}


\footnotetext[1]{Тихоокеанский государственный университет, aobelozerov@gmail.com}


\footnotetext[2]{Тихоокеанский государственный университет, mazur@khb.ru}





 \label{st\stat}





 \thispagestyle{headings}


 


 \vspace*{-12pt} 


 


  


  


\Abst{Предложен новый метод экстраполяции вариационных расчетов, в~основе которого 


лежит обуче\-ние большого чис\-ла искусственных нейронных сетей (ИНС) с~по\-сле\-ду\-ющей 


фильт\-ра\-ци\-ей для отбора сетей, обучен\-ных наилучшим образом. Метод апробирован на 


модельной задаче с~точ\-ным решением и~применен к~оценке энергии связи ядра ${}^4$He 


на основе расчетов в~модели оболочек без инертного кора (МОБИК) с~реалистическим потенциалом 


\mbox{Daejeon16}. Изучена сходимость результатов при увеличении объема входных данных. 


Показано, что метод обеспечивает достаточно высокую точ\-ность прогноза даже в~случае 


небольших модельных пространств. Предложенный метод можно применять как 


для на\-хож\-де\-ния других характеристик ядер, так и~для решения иных задач, не связанных с~ядер\-ной физикой.}





\KW{машинное обучение; методы экстраполяции; энергия основного состояния; модель 


ядерных оболочек}





\DOI{10.14357\!/\!08696527220202} 





\vskip 10pt plus 9pt minus 6pt





 \thispagestyle{headings}


 


 %\vspace*{-20pt}











\section{Введение}





  В настоящее время методы машинного обуче\-ния широко используются для 


решения различных задач. Однако у~этих методов есть два основных 


ограничения для реализации. Во-пер\-вых, как и~большинство других 


статистических методов, машинное обуче\-ние требует для корректной работы 


большого объема данных. Обычно хорошие наборы данных со\-сто\-ят минимум 


из 1000~примеров. Второе ограничение относится к~типу проб\-ле\-мы. Обычно 


машинное обуче\-ние используется в~задачах интерполяции, в~то время как его 


использование в~задачах экстраполяции изучено недостаточно.


  


  Модель оболочек без инертного кора~[1] в~на\-сто\-ящее время 


является одним из основных методов описания атом\-ных ядер. Этот подход 


называется \textit{ab initio}, поскольку в~качестве вход\-ной информации 


используется только реалистичное нук\-лон-нук\-лон\-ное взаимодействие 


(различные типы реалистичных потенциалов пред\-став\-ле\-ны в~[2--5]). 


  


  Результаты расчетов в~оболочечных моделях, вклю\-чая \mbox{МОБИК}, зависят от 


двух па\-ра\-мет\-ров: размера базисного пространства модели, опре\-де\-ля\-емо\-го 


максимальным чис\-лом осцилляторных квантов $N_{\max}$, и~величины 


осцилляторной энергии $\hbar \Omega$.


  


  


  Основная проблема \mbox{МОБИК}~--- резкое увеличение объема тре\-бу\-емых 


вы\-чис\-ли\-тель\-ных ресурсов с~увеличением $N_{\max}$, что связано 


с~экспоненциальным рос\-том чис\-ла базисных функций. Как следствие, 


современные суперкомпьютеры поз\-во\-ля\-ют проводить рас\-че\-ты в~\mbox{МОБИК} 


с~$N_{\max}\hm\leq 20$ только для самых лег\-ких ядер. 


  


  Более точные результаты могут быть получены с~использованием прос\-тых 


методов экстраполяции результатов \mbox{МОБИК} на бесконечно большие 


модельные пространства, таких как экстраполяция~B~[6, 7]. Было разработано 


множество методов для экстраполяции вы\-чис\-ле\-ний в~более слож\-ных ядерных 


сис\-те\-мах~[8--13]. Однако все они не имеют строгого обоснования, а~потому 


вопрос поиска новых методов экстраполяции остается актуальным.


  


  Цель данной работы~--- разработка нового метода экстраполяции 


результатов вариационных вы\-чис\-ле\-ний на бесконечные модельные 


пространства с~использованием методов машинного обуче\-ния, способного 


работать с~малыми наборами данных. Полученный алгоритм протестирован на 


модельной задаче с~известным точ\-ным решением и~применен к~задаче 


экстраполяции энергии основного со\-сто\-яния~${}^4$He на случай бесконечно 


больших модельных пространств.


  


\section{Модельная задача}





  Рассмотрим движение квантовой частицы с~приведенной массой $\mu 


c^2\hm= 751$~МэВ в~поле сфе\-ри\-чески-сим\-мет\-рич\-но\-го потенциала 


в~парциальной волне $l\hm=1$ с~пол\-ным угловым моментом $j\hm=1/2$. 


Радиальное урав\-не\-ние Шредингера имеет вид:


  \begin{equation}


  H u_l(k,r)=E u_l(k,r)\,.


  \label{e1-b}


  \end{equation}


Здесь гамильтониан $H\hm= T\hm+ V$, где $T$~--- оператор кинетической энергии. 


В~качестве модельного потенциала~$V$ используется потенциал  


Вуд\-са--Сак\-со\-на с~поверхностными спин-ор\-би\-таль\-ны\-ми силами:


\begin{equation*}


V(r)=  \fr{V_0}{1+\exp ((r-R)/a)} + \fr{2(\mathbf{l}\cdot 


\mathbf{s})V_{ls}}{ra}\,\fr{\exp ((r-R)/a)}{(1+\exp ((r-R)/a))^2}\,,


%\label{e2-b}


\end{equation*}


где $V_0=-30$~МэВ; $V_{ls}\hm= -10$~МэВ$\cdot$фм$^2$; $R\hm= 


3{,}08$~фм; $a\hm= 0{,}53$~фм.


  


  В осцилляторном базисе уравнение~(1) эквивалентно бесконечной сис\-те\-ме 


линейных ал\-геб\-ра\-и\-че\-ских уравнений на волновые функ\-ции в~осцилляторном 


пред\-став\-ле\-нии~$a_{Nl}(k)$:


  \begin{equation}


  \sum\limits_{N^\prime =N_0, N_0+2, \ldots , \infty} \hspace*{-8mm}a_{N^\prime l}(k) \left( 


H_{NN^\prime}-\delta_{NN^\prime}E\right) =0\,,\enskip N=N_0, N_0+2, \ldots , 


\infty\,.


  \label{e3-b}


  \end{equation}


Здесь $H_{NN^\prime}$~--- 


мат\-рич\-ные элементы гамильтониана:


$$


H_{NN^\prime} = \int\limits_0^\infty R_{Nl} HR_{N^\prime l}\, dr\,,


$$ 





  \vspace*{-3pt}


  


  \noindent


где $R_{Nl}$~--- радиальная осцилляторная 


функция; $N\hm= 2n\hm+l$~--- пол\-ное (энергетическое) чис\-ло осцилляторных 


кван\-тов, где $n$~--- глав\-ное кван\-то\-вое чис\-ло; $N_0$~--- минимальное чис\-ло 


осцилляторных кван\-тов рас\-смат\-ри\-ва\-емой сис\-те\-мы (в~данном случае  


$N_0\hm=l$).





  


  Для короткодействующих потенциалов уравнение~(\ref{e3-b}) сводится 


к~конечной сис\-те\-ме уравнений:





\vspace*{-8pt}





\noindent


  \begin{multline}


  \sum\limits_{N^\prime =N_0, N_0+2, \ldots , N_0+N_{\max}} \hspace*{-10mm} a_{N^\prime l}(k) 


\left( H_{NN^\prime} -\delta_{NN^\prime} E\right)=0\,,\\ 


  N=N_0, N_0+2, \ldots , N_0+N_{\max}\,,


  \label{e4-b}


  \end{multline}


  


  \vspace*{-6pt}


  


  \noindent


где $N_{\max}$~--- число квантов воз\-буж\-де\-ния, за\-да\-ющее модельное 


пространство.





 \begin{figure}[b] %fig1


  \vspace*{-6pt}


\begin{center}


 \mbox{%


 \epsfxsize=73.376mm 


\epsfbox{bel-1.eps}


 }


\end{center}


\vspace*{-10pt}


  \Caption{Вариационные расчеты в~модельной задаче: \textit{1}~--- $N_{\max}\hm=4$; 


\textit{2}~--- 8; \textit{3}~--- 12; \textit{4}~--- 24; \textit{5}~--- 30; \textit{6}~--- $N_{\max}\hm=38$}


  \end{figure}





  После диагонализации гамильтониана~$H$ в~(\ref{e4-b}) получаем спектр 


собственных энергий~$E_\lambda$.  Минимальное значение функции~$E_{gs}(\hbar\Omega)$ 


для данного модельного пространства ас\-со\-ци\-иру\-ет\-ся 


с~энергией основного со\-сто\-яния яд\-ра~$E_{gs}$. Отметим, что для модельной 


задачи размер мат\-ри\-цы гамильтониана с~рос\-том $N_{\max}$ увеличивается 


линейно.


  


  Метод имеет два параметра: значение осцилляторной энергии $\hbar \Omega$ 


и~максимальное чис\-ло кван\-тов воз\-буж\-де\-ния $N_{\max}$. Были проведены 


оболочечные расчеты $E_{gs}(\hbar \Omega)$ в~различных модельных 


пространствах $N_{\max}\hm= 4\ldots40$ с~шагом в~1~МэВ по $\hbar \Omega$. 


На рис.~1 показаны данные, полученные в~отдельных модельных 


пространствах.


  


 


  


  


  Горизонтальная сплош\-ная линия обозначает точ\-ное значение 


$E_{gs}^{\mathrm{exp}}\hm= -3{,}5814$~МэВ энергии основного со\-сто\-яния сис\-те\-мы, 


полученное прямым интегрированием уравнения Шредингера. Как показано 


в~\cite{14-b}, некоторые методы экстраполяции работают лучше, когда 


используются только данные, лежащие правее минимумов кривых 


$E_{gs}(\hbar\Omega)$ (вертикальная штриховая линия). Предыду\-щие работы 


авторов статьи под\-тверж\-да\-ют это утверж\-де\-ние, дополняя его фактом 


улучшения про\-из\-во\-ди\-тель\-ности после использования данных только с~$E_{gs} 


\hm < 0$~МэВ (горизонтальная штрихпунктирная линия).





\vspace*{-2pt}





\section{Задача нахождения энергии основного состояния ядра ${}^4$He}





\vspace*{-2pt}








  В МОБИК волновая функция разложена по мно\-го\-час\-тич\-но\-му 


осцилляторному базису. Как и~в~случае описанной выше модельной задачи, 


энергия основного со\-сто\-яния~$E_{gs}$ является функцией максимального 


чис\-ла осцилляторных кван\-тов $N_{\max}$ и~па\-ра\-мет\-ра осцилляторной 


функ\-ции $\hbar \Omega$. Но, в~отличие от модельной задачи, раз\-мер 


мно\-го\-час\-тич\-но\-го базиса с~увеличением $N_{\max}$ увеличивается не линейно, 


а~экспоненциально.





\begin{figure}[b] %fig2


\vspace*{-6pt}


\begin{center}


 \mbox{%


 \epsfxsize=75.092mm 


\epsfbox{bel-2.eps}


 }


\end{center}


\vspace*{-10pt}


\Caption{Энергия основного со\-сто\-яния ядра ${}^4$He $E_{gs}$, вы\-чис\-лен\-ная в~\mbox{МОБИК} 


с~реалистическим взаимодействием Daejeon16: \textit{1}~--- $N_{\max}\hm= 2$; \textit{2}~--- 4; 


\textit{3}~--- 8; \textit{4}~--- 14; \textit{5}~--- $N_{\max}\hm=18$}


\end{figure}


  


  На рис.~2 показаны результаты расчетов энергии $E_{gs}$ основного 


со\-сто\-яния ядра атома гелия ${}^4$He в~\mbox{МОБИК} с~ре\-а\-ли\-сти\-че\-ским 


взаимодействием Daejeon16~\cite{5-b} для некоторых значений $N_{\max}$. 


Расчеты проводились в~модельных пространствах $N_{\max}\hm= 2\ldots18$ при 


$\hbar\Omega\hm= 8$, 9 и~10~МэВ и~далее с~шагом 2,5~МэВ до 50~МэВ. Как 


и~в~модельной задаче, данные, лежащие слева от минимумов $E(\hbar\Omega)$ 


(вертикальная штриховая линия), не использовались в~процессе обуче\-ния 


ИНС.


  











  Точное значение энергии основного состояния может быть получено 


в~пределе $N_{\max}\hm\to \infty$. Однако из-за экспоненциального рос\-та 


чис\-ла базисных функций вы\-чис\-ле\-ние собственных значений мат\-ри\-цы 


гамильтониана для больших $N_{\max}$ становится не\-воз\-мож\-ным даже на 


суперкомпьютерах. Поэтому актуальна проб\-ле\-ма экстраполяции результатов, 


полученных в~конечных модельных пространствах, на случай $N_{\max}\hm\to 


\infty$.


  


  Оценка энергии основного состояния, полученная в~\cite{5-b} путем 


экстраполяции результатов \mbox{МОБИК} на случай бесконечно больших модельных 


пространств, со\-став\-ля\-ет $-28{,}372$~МэВ; экспериментальное значение 


энергии связанного со\-сто\-яния~--- $-28{,}296$~МэВ.





\vspace*{-2pt}


  


\section{Метод экстраполяции}





\vspace*{-2pt}














  Как уже упоминалось ранее, методы машинного обуче\-ния обычно не 


используются для задач экстраполяции с~малым объемом данных. Тем не менее 


на их основе может быть построен алгоритм для их решения. Оказывается, для 


таких задач подходит одновременное обуче\-ние множества ИНС.


  


  Предсказания, полученные с~помощью методов машинного обуче\-ния, 


включают некоторую слу\-чай\-ность, связанную с~начальными значениями 


весовых коэффициентов. Чтобы избежать этой слу\-чай\-ности, обуча\-ет\-ся большое 


чис\-ло ИНС с~одинаковыми ги\-пер\-па\-ра\-мет\-ра\-ми, но с~разными начальными 


значениями весовых коэффициентов. 








  


  Очевидно, что не все обученные ИНС будут давать правильные прогнозы, 


поэтому важ\-но иметь критерии, поз\-во\-ля\-ющие от\-фильт\-ро\-вы\-вать заведомо 


неправильно обучен\-ные ИНС.


  


  В модели ядерных оболочек роль такого критерия играет вариационный 


принцип для энергий связанных со\-сто\-яний, со\-глас\-но которому энергия 


основного\linebreak\vspace*{-12pt}





\begin{floatingfigure}{70mm} %fig3


  \vspace*{6pt}


\begin{center}


 \mbox{%


 \epsfxsize=62.058mm 


\epsfbox{bel-3.eps}


 }


\end{center}


\vspace*{-10pt}


  \Caption{Гистограмма зависимости числа ИНС, пред\-ска\-зы\-ва\-ющих энергию в~каж\-дом 


интервале, от пред\-ска\-зы\-ва\-емых ими энергий основного со\-сто\-яния~$E_{{gs}}$:


\textit{1}~--- $E_{{gs}}^{\mathrm{pr}}\hm= -3{,}564$~МэВ; $\sigma\hm= 0{,}0029$~МэВ}


\vspace*{-24pt}


\end{floatingfigure}





\noindent


 со\-сто\-яния~$E_{gs}$ долж\-на уменьшаться с~рос\-том~$N_{\max}$ при 


каж\-дом фиксированном значении~$\hbar\Omega$. 


  


  Для дальнейшей фильтрации, аналогично работе~\cite{13-b}, выбираются 


50~ИНС с~наименьшим сред\-не\-квад\-ра\-тич\-ным отклонением предсказанных 


значений от истинных.








  


  Эмпирически было показано, что за\-ви\-си\-мость чис\-ла ИНС, пред\-ска\-зы\-ва\-ющих 


энергию в~каждом интервале, от пред\-ска\-зы\-ва\-емых ими энергий основного 


со\-сто\-яния~$E_{\mathrm{gs}}$ в~большинстве случаев пред\-став\-ля\-ет собой нормальное 


распределение. Аппроксимируя их кривой Гаусса, мож\-но получить\hspace*{74.5mm}\linebreak


 прогноз 


метода и~его по\-греш-\hspace*{74.5mm}\linebreak


ность. Среднее значение этого распределения используется 


как значение, пред\-ска\-зы\-ва\-емое методом,~--- $E_{gs}^{\mathrm{pr}}$, а~стан\-дарт\-ное 


отклонение распределения~$\sigma$ принимается за по\-греш\-ность~$\Delta E$ 


пред\-ска\-за\-ния метода.


  %


  На рис.~3 приведена такая гис\-то\-грам\-ма. Гаус\-со\-и\-да для данного 


распределения показана кривой. 


%Прогнозируемое значение 


%энергии~$E_{gs}^{pr}$ и~стан\-дарт\-ное отклонение~$\sigma$ показаны в~правом 


%верх\-нем углу рисунка.


  





  


  





  Вышеупомянутый алгоритм был разработан с~использованием биб\-лио\-те\-ки 


с~открытым исходным кодом Keras~\cite{15-b} для языка Python. 





\vspace*{-2pt}


  


\section{Результаты}





\vspace*{-2pt}





  С помощью поиска по сетке были получены оптимальные па\-ра\-мет\-ры 


машинного обуче\-ния для обеих задач. К~этим па\-ра\-мет\-рам относятся: ско\-рость 


обуче\-ния $\gamma\hm= 0{,}1$, размер пакета $B_S\hm=64$ и~16 для модельной 


задачи и~задачи на\-хож\-де\-ния энергии основного со\-сто\-яния ядра ${}^4$He 


соответственно. Оптимальная архитектура ИНС в~обоих случаях~--- 5~скрытых 


слоев по~10~нейронов в~каж\-дом. Ак\-ти\-ви\-ру\-ющей функцией каж\-до\-го из слоев, 


за исключением по\-след\-не\-го, является сиг\-мо\-и\-даль\-ная функция.





\begin{figure}[b] %fig4


\vspace*{-6pt}


\begin{center}


 \mbox{%


 \epsfxsize=126.9mm 


\epsfbox{bel-4.eps}


 }


\end{center}


\vspace*{-10pt}


\Caption{Результаты работы алгоритма с~разным объемом входных данных для модельной 


задачи~(\textit{а}) и~для задачи 


на\-хож\-де\-ния энергии основного со\-сто\-яния ${}^4$He~(\textit{б}); $\gamma\hm=0{,}1$}


\end{figure}





  Зависимость предсказаний метода от объема данных для модельной задачи 


показана на рис.~4,\,\textit{а}. Точки пред\-став\-ля\-ют собой пред\-ска\-зан\-ное значение энергии 


основного со\-сто\-яния $E_{gs}$, вертикальные линии~--- по\-греш\-ность 


предсказания~$\sigma$. Горизонтальная линия соответствует истинному 


значению энергии $E_{gs}^{\mathrm{exp}}\hm= -3{,}581$~МэВ. Процесс обучения 


проводился на данных вплоть до $N_{\max}\hm= 10$, 12, 14 и~так далее. Как 


видно из рис.~4,\,\textit{а}, истинное значение лежит в~пред\-ска\-зы\-ва\-емых интервалах для 


всех размеров набора данных. При увеличении максимального значения 


$N_{\max}$ пред\-ска\-за\-ния метода изменяются с~$-3{,}585\hm\pm 0{,}030$ 


до $-3{,}579\hm\pm 0{,}006$~МэВ для $N_{\max}\hm=10$ и~40 соответственно.


  


  








  На рис.~4,\,\textit{б} показана аналогичная за\-ви\-си\-мость для задачи на\-хож\-де\-ния энергии 


основного со\-сто\-яния ${}^4$He. Горизонтальная линия~--- результаты 


экстраполяции, полученные в~\cite{5-b}. Все предсказанные диапазоны так\-же 


содержат это значение. Пред\-ска\-за\-ния для $N_{\max}\hm=10$ 


и~$18$ со\-став\-ля\-ют $-28{,}368\hm\pm 0{,}043$  


и~$-28{,}366 \hm\pm 0{,}015$~МэВ соответственно.


  


   


  


  Разработанный алгоритм хорошо работает даже с~небольшими наборами 


данных. Так, задача на\-хож\-де\-ния энергии основного со\-сто\-яния~${}^4$He 


с~$N_{\max}\hm=10$ содержит около~50~примеров, что является 


недостаточным объемом данных для корректной работы большинства других 


алгоритмов на основе машинного обуче\-ния.


  


\section{Заключение}





  Предложен новый метод экстраполяции результатов вариационных расчетов 


  в~модели ядерных оболочек с~осцилляторным базисом на случай бесконечно 


больших модельных пространств. Метод основан на машинном обуче\-нии 


и~поз\-во\-ля\-ет пред\-ска\-зы\-вать значения энергии основного со\-сто\-яния на основе 


результатов вариационных расчетов, полученных в~ограниченном чис\-ле 


модельных пространств $N_{\max}$ с~различными значениями осцилляторного 


па\-ра\-мет\-ра~$\hbar\Omega$. В~разработанном алгоритме используется только 


программное обеспечение с~открытым исходным кодом.


  


  В рамках модельной задачи с~известным точ\-ным решением исследовалась 


схо\-ди\-мость результатов при увеличении объема входных данных. Показано, 


что метод поз\-во\-ля\-ет до\-стичь достаточно высокой точ\-ности прогноза даже 


в~случае малых пространств моделей $N_{\max}\sim 10$, что особенно важ\-но 


с~точ\-ки зрения его применения для анализа результатов, полученных 


в~современных оболочечных подходах. То же утверж\-де\-ние справедливо и~для 


оцен\-ки энергии основного со\-сто\-яния ядра ${}^4$He.


  


  Метод позволяет не только получить ин\-те\-ре\-су\-ющее значение, но и~оценить 


точ\-ность прогнозов. К~достоинствам метода мож\-но так\-же отнести 


уни\-вер\-саль\-ность: его можно использовать не только для анализа энергий 


основного со\-сто\-яния, но и~для других характеристик ядер.





{\small\frenchspacing


{%\baselineskip=10.8pt


\addcontentsline{toc}{section}{Литература}


\begin{thebibliography}{99}


\bibitem{1-b}


\Au{Barrett B.\,R., Navr\'{a}til~P., Vary~J.\,P.} Ab initio no core shell model~/\!\!/ 


Prog. Part. Nucl. Phys., 2013. Vol.~69. Iss.~1. P.~131--181. doi: 10.1016\!/\!j.ppnp. 2012.10.003.





\bibitem{4-b} %2


\Au{Stoks V.\,G.\,J., Klomp R.\,A.\,M., Terheggen~C.\,P.\,F., de Swart~J.\,J.} Construction of high-quality 


NN potential models~/\!\!/ Phys. Rev.~C, 1994. Vol.~49. P.~2950--2962. doi: 


10.1103\!/\!PhysRevC.49.2950.











\bibitem{3-b} %3


\Au{Wiringa R.\,B., Stoks~V.\,G.\,J., Schiavilla~R.} Accurate nucleon--nucleon potential with  


charge-independence breaking~/\!\!/ Phys. Rev.~C,  1995. Vol.~51. P.~38--51. doi: 


10.1103\!/\!PhysRevC.51.38.





\bibitem{2-b} %4


\Au{Machleidt R.} High-precision, charge-dependent Bonn nucleon--nucleon potential~/\!\!/ Phys. Rev.~C, 


2001.  Vol.~63. Iss.~2. Art.~024001. 32~p. doi: 10.1103\!/\!PhysRevC.63. 024001.








\bibitem{5-b}


\Au{Shirokov A.\,M., Shin I.\,J., Kim~Y., Sosonkina~M., Maris~P., Vary~J.\,P.} N3LO NN interaction 


adjusted to light nuclei in ab exitu approach~/\!\!/ Phys. Lett.~B, 2016. Vol.~761. P.~87--91. doi: 


10.1016\!/\!j.physletb.2016.08.006.


\bibitem{6-b}


\Au{Maris~P., Vary~J.\,P., Shirokov~A.\,M.} Ab initio no-core full configuration calculations of light 


nuclei~/\!\!/ Phys. Rev.~C, 2009. Vol.~79. Art.~014308. 15~p. doi: 10.1103\!/\!PhysRevC. 79.014308. 


\bibitem{7-b}


\Au{Maris P., Vary J.\,P.} Ab initio nuclear structure calculations of p-shell nuclei with JISP16~/\!\!/ Int. 


J.~Mod. Phys.~E, 2013. Vol.~22. Iss.~7. Art. 1330016. 17~p. doi: 10.1142\!/\!S0218301313300166.














\bibitem{10-b} %8


\Au{Zhan H., Nogga~A., Barrett B.\,R., Vary~J.\,P., Navr\'{a}til~P.} Extrapolation method for the no-core 


shell model~/\!\!/ Phys. Rev.~C, 2004. Vol.~69. Art. 034302. 8~p. doi: 10.1103\!/ PhysRevC.69.034302. 





\bibitem{8-b} %9


\Au{Coon S.\,A., Avetian M.\,I., Kruse M.\,K.\,G., van Kolck~U., Maris~P., Vary~J.\,P.} Convergence 


properties of ab initio calculations of light nuclei in a harmonic oscillator basis~/\!\!/ Phys. Rev.~C, 2012. 


Vol.~86. Art. 054002. doi: 10.1103\!/\!PhysRevC.86.054002.








\bibitem{12-b} %10


\Au{Kruse M.\,K.\,G., Jurgenson~E.\,D., Navr\'{a}til~P., Barrett~B.\,R., Ormand~W.\,E.} Extrapolation 


uncertainties in the importance-truncated no-core shell model~/\!\!/ Phys. Rev.~C, 2013. Vol.~87. 


Art. 044301. 15~p. doi: 10.1103\!/\!PhysRevC.87.044301.





\bibitem{11-b}


\Au{More S.\,N., Ekstr\"{o}m~A., Furnstahl R.\,J., Hagen~G., Papenbrock~T.} Universal properties of 


infrared oscillator basis extrapolations~/\!\!/ Phys. Rev.~C, 2013. Vol.~87. Art. 044326. 14~p. doi: 


10.1103\!/\!PhysRevC.87.044326.





\bibitem{9-b} %12


\Au{Furnstahl R.\,J., More~S.\,N., Papenbrock~T.} Systematic expansion for infrared oscillator basis 


extrapolations~/\!\!/ Phys. Rev.~C, 2014.  Vol.~89. Art. 044301. 12~p. doi: 10.1103\!/ PhysRevC.89.044301. 





\bibitem{13-b}


\Au{Negoita G.\,A., Vary~J.\,P., Luecke~G.\,R., Maris~P., Shirokov~A.\,M., Shin~I.\,J., Kim~Y., Ng~E.\,G.,  


Yang~C., Lockner~M., Prabhu~G.\,M.} Deep learning: Extrapolation tool for ab initio nuclear theory~/\!\!/ 


Phys. Rev.~C, 2019. Vol.~99. Art. 054308. 11~p. doi: 10.1103\!/\!PhysRevC.99.054308. 


\bibitem{14-b}


\Au{Shirokov A.\,M., Mazur A.\,I., Kulikov~V.\,A.} On the convergence of oscillator basis calculations~/\!\!/ 


Phys. Atom. Nucl., 2021. Vol.~84. Iss.~2. P.~131--143. doi: 10.1134\!/ S1063778821020149.


\bibitem{15-b}


Keras: The Python Deep Learning library. {\sf https://keras.io}.


       \end{thebibliography}


} }











\vspace*{-3pt}





\hfill{\small\textit{Поступила в~редакцию 06.12.21}}





\vspace*{8pt}





\hrule





\vspace*{2pt}





%\newpage





%\vspace*{-28pt}





\hrule





%\vspace*{8pt}





\def\tit{MACHINE LEARNING FOR~EXTRAPOLATION PROBLEMS WITH~SMALL DATASET}











\def\titkol{Machine learning for~extrapolation problems with~small dataset}





\def\aut{A.\,O.~Belozerov and~A.\,I.~Mazur}





\def\autkol{A.\,O.~Belozerov and~A.\,I.~Mazur}





\titel{\tit}{\aut}{\autkol}{\titkol}





\vspace*{-8pt}





\noindent


Pacific National University, 136~Tikhookeanskaya Str., Khabarovsk 680035, Russian Federation





\def\leftfootline{\small{\textbf{\thepage}


\hfill Sistemy i Sredstva Informatiki~---


Systems and Means of Informatics 2022 vol~32 no\,2}


}%


 \def\rightfootline{\small{Sistemy i Sredstva Informatiki~---


 Systems and Means of Informatics 2022 vol~32 no\,2


\hfill \textbf{\thepage}}}





\vspace*{3pt}











\Abste{A new method of extrapolation of variational calculations which is based on 


training a~large number of artificial neural networks with subsequent


filtering to select the best 


networks  is proposed. This method is used to evaluate\linebreak\vspace*{-12pt}}





\Abstend{ground state energy for model problem and to calculate 


the ground state energy of a~${}^{4}$He nucleus based on calculations in no-core shell model with 


realistic potential Daejeon16. The results convergence with increasing number of input data was 


studied. It is shown that the method allows one to achieve a~sufficiently high prediction accuracy 


even in the case of small model spaces. The method can be applied both in finding various 


characteristics of nuclei and in solving other problems not related to nuclear physics.}





\KWE{machine learning; extrapolation methods; ground state energy; nuclear shell model }














\DOI{10.14357\!/\!08696527220202} 





%\vspace*{-16pt}








 \Ack


\noindent


The paper is published on the proposal of the Program Committee of the 6th International Scientific and 


Practical Conference ``Information Technologies and High Performance Computing'' (ITHPC-2021). The 


work was partially supported by the Russian Foundation for Basic Research (project No.\,20-02-00357) and by 


the Ministry of Science and Higher Education of the Russian Federation (project No.\,0818-2020-0005).


 








%\vspace*{-6pt}





\renewcommand{\bibname}{\protect\rmfamily References}


%\renewcommand{\bibname}{\large\protect\rm References}





{\small\frenchspacing


{%\baselineskip=10.8pt


\addcontentsline{toc}{section}{References}


\begin{thebibliography}{99}


\bibitem{1-b-1}


\Aue{Barrett, B.\,R., P.~Navr\'{a}til, and J.\,P.~Vary.} 2013. Ab initio no core shell model. \textit{Prog. 


Part. Nucl. Phys.} 69(1):131--181. doi: 10.1016\!/\!j.ppnp.2012.10.003.





\bibitem{4-b-1} %2


\Aue{Stoks, V.\,G.\,J., R.\,A.\,M.~Klomp, C.\,P.\,F.~Terheggen, and J.\,J.~de Swart.} 1994. Construction of  


high-quality NN potential models. \textit{Phys. Rev.~C} 49:2950--2962. doi: 


10.1103\!/\!PhysRevC.49.2950.











\bibitem{3-b-1}


\Aue{Wiringa, R.\,B., V.\,G.\,J.~Stoks, and R.~Schiavilla.} 1995. Accurate nucleon--nucleon potential with 


charge-independence breaking. \textit{Phys. Rev.~C} 51:38--51. doi: 10.1103\!/ PhysRevC.51.38.








\bibitem{2-b-1} %4


\Aue{Machleidt, R.} 2001. High-precision, charge-dependent Bonn nucleon--nucleon potential. \textit{Phys. 


Rev.~C} 63(2):024001. 32~p. doi: 10.1103\!/\!PhysRevC.63.024001.





\bibitem{5-b-1} %5


\Aue{Shirokov, A.\,M., I.\,J.~Shin, Y.~Kim, M.~Sosonkina, P.~Maris, and J.\,P.~Vary.} 2016. N3LO  


NN interaction adjusted to light nuclei in ab exitu approach. \textit{Phys. Lett.~B} 761:87--91. doi: 


10.1016\!/\!j.physletb.2016.08.006.


\bibitem{6-b-1}


\Aue{Maris, P., J.\,P.~Vary, and A.\,M.~Shirokov.} 2009. Ab initio no-core full configuration calculations 


of light nuclei. \textit{Phys. Rev.~C} 79:014308. 15~p. doi: 10.1103\!/ PhysRevC.79.014308.


\bibitem{7-b-1}


\Aue{Maris, P., and J.\,P.~Vary.} 2013. Ab initio nuclear structure calculations of \mbox{p-shell} nuclei with 


JISP16. \textit{Int. J.~Mod. Phys.~E} 22(7):1330016. 17~p. doi: 10.1142\!/ S0218301313300166.








\bibitem{10-b-1} %8


\Aue{Zhan, H., A.~Nogga, B.\,R.~Barrett, J.\,P.~Vary, and P.~Navr\'{a}til.} 2004. Extrapolation method for 


the no-core shell model. \textit{Phys. Rev.~C} 69:034302. 8~p. doi: 10.1103\!/ PhysRevC.69.034302.





\bibitem{8-b-1} %9


\Aue{Coon, S.\,A., M.\,I.~Avetian, M.\,K.\,G.~Kruse, U.~van Kolck, P.~Maris, and J.\,P.~Vary.} 2012. 


Convergence properties of ab initio calculations of light nuclei in a~harmonic oscillator basis. \textit{Phys. 


Rev.~C} 86:054002. 15~p. doi: 10.1103\!/\!PhysRevC.86.054002.





\bibitem{12-b-1} %10


\Aue{Kruse, M.\,K.\,G., E.\,D.~Jurgenson, P.~Navr\'{a}til, B.\,R.~Barrett, and W.\,E.~Ormand.} 2013. 


Extrapolation uncertainties in the importance-truncated no-core shell model. \textit{Phys. Rev.~C} 


87:044301. 15~p. doi: 10.1103\!/\!PhysRevC.87.044301.





\bibitem{11-b-1}


\Aue{More, S.\,N., A.~Ekstr\"{o}m, R.\,J.~Furnstahl, G.~Hagen, and T.~Papenbrock.} 2013. Universal 


properties of infrared oscillator basis extrapolations. \textit{Phys. Rev.~C} 87:044326. 14~p. doi: 


10.1103\!/\!PhysRevC.87.044326.





\bibitem{9-b-1} %12


\Aue{Furnstahl, R.\,J., S.\,N.~More, and T.~Papenbrock.} 2014. Systematic expansion for infrared oscillator 


basis extrapolations. \textit{Phys. Rev.~C} 89:044301. 12~p. doi: 10.1103\!/ PhysRevC.89.044301.








\bibitem{13-b-1}


\Aue{Negoita, G.\,A., J.\,P.~Vary, G.\,R.~Luecke, P.~Maris, A.\,M.~Shirokov, I.\,J.~Shin, Y.~Kim, 


E.\,G.~Ng, C.~Yang, M.~Lockner, and G.\,M.~Prabhu.} 2019. Deep learning: Extrapolation tool for ab 


initio nuclear theory. \textit{Phys. Rev.~C} 99:054308. 11~p. doi: 10.1103\!/\!PhysRevC.99.054308.


\bibitem{14-b-1}


\Aue{Shirokov, A.\,M., A.\,I.~Mazur, and V.\,A.~Kulikov.} 2021. On the convergence of oscillator basis 


calculations. \textit{Phys. Atom. Nucl.} 84(2):131--143. doi: 10.1134\!/ S1063778821020149.


\bibitem{15-b-1}


Keras: The Python deep learning library. Available at: {\sf https://keras.io/} (accessed March~28, 2022).


\end{thebibliography}


} }





\vspace*{-3pt}





\hfill{\small\textit{Received December 6, 2021}} 





%\vspace*{-14pt} 





\Contr





\noindent


\textbf{Belozerov Alexander O.} (b.\ 2000)~--- laboratory assistant, Research Laboratory for Modeling 


Quantum Processes, Pacific National University, 136~Tikhookeanskaya Str., Khabarovsk 680035, Russian 


Federation; \mbox{aobelozerov@gmail.com}





\vspace*{3pt}





\noindent


\textbf{Mazur Alexander I.} (b.\ 1959)~--- leading scientist, Research Laboratory for Modeling Quantum 


Processes, Pacific National University, 136~Tikhookeanskaya Str., Khabarovsk 680035, Russian Federation; 


\mbox{mazur@khb.ru}





\label{end\stat}





\renewcommand{\bibname}{\protect\rm Литература} 


